% mode: ini-etex
% Aimed at e-TeX's \savinghyphcodes (sections 1854-1858): hyph_store_codes keeps the \lccode table per language in the trie, section 1857 packs it, and set_hyph_index and set_lc_code use it at hyphenation and in \hyphenation.  A language saved twice
% (the second table replaces the first, trie_c overwritten in section 1856), one with \lccode changed afterwards, which does not affect it; patterns only in some; "Not a letter" judged by the saved code, after the trie is ready, and by \lccode, before; c=255 being forced in when no code is set.
\input prelude
\lefthyphenmin=1 \righthyphenmin=1 \defaulthyphenchar=`\- \hyphenpenalty=-1000 \hsize=1pt \tolerance=10000 \parfillskip=0pt \parindent=0pt \pretolerance=-1
\tenrm \hyphenchar\tenrm=`\-
\def\try#1{\setbox0=\vbox{\hskip0pt #1\par}\showbox0 }
\savinghyphcodes=1
\language=1
\patterns{a1b b1c}
% Saved twice, with a change between.
\language=2
\lccode`B=`b
\patterns{a1b}
\lccode`B=0 \lccode`Z=`b
\patterns{b1c}
\lccode`Z=`z \lccode`B=`b
% Saved with no letters at all: c=255 is forced in, section 1856.
\language=3
\begingroup
\count0=0
\def\clear{\lccode\count0=0 \advance\count0 by 1 \ifnum\count0<256 \expandafter\clear\fi}
\clear
\patterns{}
\endgroup
% A language with its codes saved, and with no patterns: "\patterns{}" did it.
\language=4
\lccode`Q=`a
\patterns{}
\lccode`Q=`q
% No saving: the live \lccode is used.
\savinghyphcodes=0
\language=5
\patterns{a1b b1c}
% Saving and language 255; a language above is 0.
\savinghyphcodes=1
\language=255
\patterns{a1b}
\language=300
\lccode`Y=`a
\patterns{e1g}
\lccode`Y=`y
\savinghyphcodes=0
% Before the trie is ready, \hyphenation uses \lccode.
\language=1
\hyphenation{A-A}
% The trie gets ready with the first hyphenation.
\language=5 \try{abc}
\language=1 \try{abc} \try{ABC}
\language=2 \try{abc} \try{ABC} \try{aZc} \try{aBc}
\language=3 \try{abc}
\language=4 \try{abc} \try{QQ}
\language=255 \try{ab}
\language=0 \try{eg} \try{Yg}
\language=9 \try{abc}
% Now it is ready: \hyphenation uses the saved codes.
\lccode`A=0 \lccode`a=0
\language=1
\hyphenation{A-b a-c}
\try{Ab}
\try{ac}
\language=5
\hyphenation{A-b a-c}
\language=3
\hyphenation{a-b}
\language=9
\hyphenation{a-b}
\lccode`A=`a \lccode`a=`a
\language=2
\lccode`B=0
\hyphenation{B-c-c}
\try{Bcc}
\end
